% Original: PDF strana 24, donji slajd; odobren drugi deo s048.
\slajdid{02-s048b}
\begin{frame}[label=02-s048b]{Dvoulazna NI kola: tri ili dva čipa?}
\setlength{\parskip}{0pt}
% element: 02-s048b:e01
Zamena troulaznog NI kola: dva dvoulazna NI i invertor, \alert{\textbf{uz povećano kašnjenje}}:
\par\vspace{2mm}
{\centering\(\displaystyle\overline{XYZ}=\overline{(XY)Z}
=\overline{\overline{\overline{XY}}\,Z}.\)\par}
\vspace{2mm}
\begin{columns}[T,onlytextwidth]
\begin{column}{.48\textwidth}
% element: 02-s048b:e02
\Oznaka{Samo dvoulazna NI kola}
\begin{itemize}
\setlength{\itemsep}{1pt}\setlength{\parsep}{0pt}
\item 4 NI kola za četiri ulazne inverzije;
\item 3 NI kola za zamenu troulaznog NI;
\item još 2 NI kola za ostatak funkcije.
\end{itemize}
\(4+3+2=9\) dvoulaznih NI kola.\par
Četiri kola po čipu: \textbf{3 ista čipa}.
\end{column}
\begin{column}{.48\textwidth}
% element: 02-s048b:e03
\Oznaka{Čip sa invertorima + čip sa NI kolima}
Čip sa šest invertora:\par
4 za ulazne komplemente;\par
1 od dva preostala za međuinverziju.
\par\vspace{1.5mm}
Čip sa četiri dvoulazna NI kola:\par
\textbf{sva četiri kola za funkciju}.
\par\vspace{1.5mm}
Ukupno: \textbf{2 čipa različitih vrsta}.
\end{column}
\end{columns}
\par\vspace{2mm}
% element: 02-s048b:e04
\Rezultat{Tri ista čipa mogu doneti popust na količinu. Da li je jeftinije \textbf{1 + 1 različitih čipova ili 3 ista}? \alert{\textbf{Nema opšteg odgovora}} — zavisi od cena i uslova nabavke.}
\end{frame}
